\documentclass[10pt]{article}

% COLM 2025 Conference Style
% Options: submission (default, with line numbers), final, preprint
\usepackage[submission]{colm2025_conference}

% Additional math commands
\input{math_commands.tex}

\usepackage{subfigure}
\usepackage{mathtools}
\usepackage[capitalize,noabbrev]{cleveref}

% Custom commands
\newcommand{\R}{\mathbb{R}}
\newcommand{\E}{\mathbb{E}}
\newcommand{\ind}{\mathbf{1}}

\title{On the Equivalence of Representational and Reasoning Superposition}

% Authors hidden during submission - will be revealed in camera-ready
\author{
  Anonymous Authors
}

\begin{document}

\maketitle

\begin{abstract}
Recent work in mechanistic interpretability has identified \emph{superposition} as a fundamental phenomenon in neural networks, where models encode more features than they have dimensions. Independently, research on chain-of-continuous-thought has revealed that language models can maintain multiple reasoning traces simultaneously in latent space---a form of \emph{reasoning superposition}. We prove that these two phenomena are mathematically equivalent, arising from the same underlying optimization problem parameterized by different co-occurrence structures. Specifically, we show that both representational superposition (packing features into limited dimensions) and reasoning superposition (maintaining multiple valid reasoning traces) minimize a unified loss function that trades off capacity against interference, weighted by co-occurrence probability. We establish formal bounds on how sparsity enables superposition in both settings and prove that temporal depth (iterative refinement) can compensate for spatial compression (bottleneck dimensions), providing theoretical foundations for architectures that combine both mechanisms.
\end{abstract}

\section{Introduction}
\label{sec:intro}

Understanding how neural networks represent and process information is a central challenge in machine learning. Two recent lines of research have independently identified \emph{superposition} as a key phenomenon, but in seemingly different contexts.

\paragraph{Representational superposition.} \citet{elhage2022toy} demonstrated that neural networks can represent more features than they have dimensions by encoding features as non-orthogonal directions. When features are sparse (rarely co-activate), the interference cost of sharing directions is low, enabling efficient compression. This phenomenon has profound implications for mechanistic interpretability: if features are superposed, they cannot be read off individual neurons, complicating efforts to understand model behavior \citep{bricken2023monosemanticity, templeton2024scaling}.

\paragraph{Reasoning superposition.} \citet{hao2024training} introduced chain-of-continuous-thought (Coconut), where language models reason in continuous latent space rather than through discrete tokens. \citet{zhu2025reasoning} proved that such models can maintain multiple reasoning traces simultaneously---a ``superposition of search paths''---enabling implicit parallel exploration. This was formalized as reasoning superposition, where the latent thought vector encodes a weighted combination of valid partial solutions.

\paragraph{The question.} Despite superficial differences, both phenomena involve packing more information into limited representational capacity. Are they fundamentally the same? Or are there deep structural differences that make them distinct phenomena requiring separate analysis?

\paragraph{Our contribution.} We prove that representational and reasoning superposition are \emph{mathematically equivalent}---both arise from minimizing the same loss function, differing only in the co-occurrence structure of the objects being encoded. Specifically:

\begin{enumerate}
    \item We establish a \textbf{unified loss function} (Theorem~\ref{thm:equivalence}) that captures both phenomena, with the co-occurrence matrix $P$ as the key distinguishing parameter.

    \item We prove \textbf{capacity bounds} (Theorem~\ref{thm:capacity}) showing how sparsity enables superposition in both settings, with matching scaling behavior.

    \item We prove a \textbf{time-space equivalence} (Theorem~\ref{thm:timespace}) showing that iterative refinement can compensate for dimensional bottlenecks, providing theoretical foundations for architectures combining both mechanisms.
\end{enumerate}

These results unify two previously separate research directions and suggest that techniques developed for one setting (e.g., sparse autoencoders for representational superposition) may transfer to the other (e.g., analyzing reasoning traces in continuous thought).

\section{Preliminaries}
\label{sec:prelim}

We establish notation and review the two superposition frameworks we aim to unify.

\subsection{Notation}

Let $[n] = \{1, \ldots, n\}$ denote the first $n$ positive integers. For vectors $\mathbf{u}, \mathbf{v} \in \R^d$, we write $\langle \mathbf{u}, \mathbf{v} \rangle = \mathbf{u}^\top \mathbf{v}$ for the inner product and $\|\mathbf{u}\| = \sqrt{\langle \mathbf{u}, \mathbf{u} \rangle}$ for the Euclidean norm. For a matrix $\mathbf{W} \in \R^{d \times n}$, we write $\mathbf{w}_i \in \R^d$ for its $i$-th column. The indicator function $\ind[\cdot]$ equals 1 if its argument is true and 0 otherwise. Throughout, $d$ denotes the \emph{bottleneck dimension}---the representational capacity constraint that induces superposition.

\subsection{Representational Superposition}
\label{sec:rep_superposition}

Following \citet{elhage2022toy}, consider an autoencoder with encoder $\mathbf{W} \in \R^{d \times n}$ and tied decoder $\mathbf{W}^\top$, where $d < n$. The model receives sparse binary inputs $\mathbf{x} \in \{0, 1\}^n$ where each component $x_i$ activates independently with probability $p_i$. The reconstruction is:
\begin{equation}
\hat{\mathbf{x}} = \sigma(\mathbf{W}^\top \mathbf{W} \mathbf{x})
\end{equation}
where $\sigma$ is a nonlinearity (typically ReLU). The model is trained to minimize importance-weighted reconstruction error:
\begin{equation}
\mathcal{L}_{\text{rep}} = \E\left[\sum_{i=1}^n \omega_i (x_i - \hat{x}_i)^2\right]
\label{eq:rep_loss}
\end{equation}
where $\omega_i \geq 0$ are importance weights.

\begin{definition}[Representational Superposition]
\label{def:rep_super}
We say features $i$ and $j$ are in \emph{superposition} if their encodings are non-orthogonal: $\langle \mathbf{w}_i, \mathbf{w}_j \rangle \neq 0$. A model exhibits superposition if it represents $n > d$ features with some pairs in superposition.
\end{definition}

The key insight of \citet{elhage2022toy} is that sparsity enables superposition: when $p_i$ are small, features rarely co-activate, so interference from non-orthogonality has low expected cost.

\subsection{Reasoning Superposition}
\label{sec:reason_superposition}

Following \citet{hao2024training} and \citet{zhu2025reasoning}, consider a transformer with continuous thought. Instead of generating discrete tokens, the model maintains a latent thought vector $\mathbf{h}_t \in \R^d$ that evolves over $T$ steps:
\begin{equation}
\mathbf{h}_{t+1} = f_\theta(\mathbf{h}_t, \mathbf{c})
\end{equation}
where $\mathbf{c}$ is the context and $f_\theta$ is a learned transformation (typically attention followed by an MLP).

For graph reachability problems, \citet{zhu2025reasoning} proved that the thought vector can be decomposed as:
\begin{equation}
\mathbf{h}_t = \sum_{k=1}^K \alpha_k^{(t)} \mathbf{v}_k
\label{eq:thought_decomp}
\end{equation}
where $\mathbf{v}_k \in \R^d$ encodes the $k$-th valid partial path and $\alpha_k^{(t)} = \text{softmax}(\text{logit}_k^{(t)})$ are attention-derived weights.

\begin{definition}[Reasoning Superposition]
\label{def:reason_super}
We say reasoning traces $i$ and $j$ are in \emph{superposition} if both contribute to the thought vector with non-zero weight: $\alpha_i^{(t)} > 0$ and $\alpha_j^{(t)} > 0$. A model exhibits reasoning superposition if it maintains $K > 1$ traces in superposition during inference.
\end{definition}

The key insight of \citet{zhu2025reasoning} is that reasoning superposition enables implicit breadth-first search: instead of committing to a single path (depth-first), the model explores multiple paths simultaneously.

\subsection{Apparent Differences}

At first glance, these phenomena seem distinct:

\begin{center}
\begin{tabular}{lcc}
\toprule
& \textbf{Representational} & \textbf{Reasoning} \\
\midrule
Objects & Input features & Reasoning traces \\
Encoding & Weight vectors $\mathbf{w}_i$ & Trace embeddings $\mathbf{v}_k$ \\
Combination & Sparse binary $\mathbf{x}$ & Soft weights $\alpha_k$ \\
Sparsity source & Feature activation & Trace validity \\
\bottomrule
\end{tabular}
\end{center}

We will show these differences are superficial: both phenomena arise from the same mathematical structure.

\section{The Equivalence Theorem}
\label{sec:equivalence}

We now prove that representational and reasoning superposition are instances of the same optimization problem.

\subsection{Unified Loss Function}

\begin{definition}[Co-occurrence Matrix]
\label{def:cooccur}
A \emph{co-occurrence matrix} $P \in [0,1]^{n \times n}$ specifies the probability that objects $i$ and $j$ are simultaneously relevant. We require $P_{ii} = p_i$ (marginal probability) and $P_{ij} = P_{ji}$ (symmetry).
\end{definition}

\begin{definition}[Superposition Loss]
\label{def:super_loss}
For encoding matrix $\mathbf{W} \in \R^{d \times n}$ (where $d$ is the bottleneck dimension) and co-occurrence matrix $P$, define:
\begin{align}
\mathcal{L}_{\text{cap}}(\mathbf{W}) &= \sum_{i=1}^{n} p_i \left(1 - \|\mathbf{w}_i\|^2\right)^2 \label{eq:cap_loss}\\
\mathcal{L}_{\text{int}}(\mathbf{W}, P) &= \sum_{i \neq j} P_{ij} \langle \mathbf{w}_i, \mathbf{w}_j \rangle^2 \label{eq:int_loss}\\
\mathcal{L}(\mathbf{W}, P) &= \mathcal{L}_{\text{cap}}(\mathbf{W}) + \mathcal{L}_{\text{int}}(\mathbf{W}, P) \label{eq:unified_loss}
\end{align}
We call $\mathcal{L}_{\text{cap}}$ the \emph{capacity loss} (penalizing deviation from unit norm) and $\mathcal{L}_{\text{int}}$ the \emph{interference loss} (penalizing non-orthogonality weighted by co-occurrence).
\end{definition}

\begin{theorem}[Representational-Reasoning Equivalence]
\label{thm:equivalence}
Both representational and reasoning superposition minimize the unified loss $\mathcal{L}(\mathbf{W}, P)$ with different co-occurrence matrices:
\begin{enumerate}
    \item[(a)] \textbf{Representational:} Setting $P_{ij} = p_i p_j$ (independent activation), minimizing $\mathcal{L}$ is equivalent to minimizing the expected reconstruction error of \citet{elhage2022toy}.
    \item[(b)] \textbf{Reasoning:} Setting $P_{ij} = \Pr[\text{traces } i, j \text{ both valid}]$, minimizing $\mathcal{L}$ is equivalent to the implicit objective of attention in continuous thought.
\end{enumerate}
\end{theorem}

\subsection{Proof of Part (a): Representational Superposition}

\begin{lemma}
\label{lem:rep_expansion}
For the autoencoder of Section~\ref{sec:rep_superposition} with ReLU activation, under the assumption that interference is small (valid when features are sparse), the expected reconstruction error satisfies:
\begin{equation}
\E\left[\|\mathbf{x} - \hat{\mathbf{x}}\|^2\right] = \sum_{i=1}^n p_i (1 - \|\mathbf{w}_i\|^2)^2 + \sum_{i \neq j} p_i p_j \langle \mathbf{w}_i, \mathbf{w}_j \rangle^2 + O(\epsilon^3)
\end{equation}
where $\epsilon = \max_{i \neq j} |p_i p_j \langle \mathbf{w}_i, \mathbf{w}_j \rangle|$ is the maximum expected interference.
\end{lemma}

\begin{proof}
For input $\mathbf{x}$ with $x_i \in \{0,1\}$ independent Bernoulli$(p_i)$, the pre-activation reconstruction of component $i$ is:
\begin{equation}
z_i = (\mathbf{W}^\top \mathbf{W} \mathbf{x})_i = \sum_{j=1}^n x_j \langle \mathbf{w}_i, \mathbf{w}_j \rangle = x_i \|\mathbf{w}_i\|^2 + \sum_{j \neq i} x_j \langle \mathbf{w}_i, \mathbf{w}_j \rangle
\end{equation}

After ReLU, $\hat{x}_i = \max(0, z_i)$. When $x_i = 1$:
\begin{equation}
z_i = \|\mathbf{w}_i\|^2 + \underbrace{\sum_{j \neq i} x_j \langle \mathbf{w}_i, \mathbf{w}_j \rangle}_{\text{interference term}}
\end{equation}

For sparse features, the interference term is typically small. Expanding to second order:
\begin{align}
\E[(x_i - \hat{x}_i)^2 | x_i = 1] &\approx \E\left[\left(1 - \|\mathbf{w}_i\|^2 - \sum_{j \neq i} x_j \langle \mathbf{w}_i, \mathbf{w}_j \rangle\right)^2\right] \\
&= (1 - \|\mathbf{w}_i\|^2)^2 + \sum_{j \neq i} p_j \langle \mathbf{w}_i, \mathbf{w}_j \rangle^2 + O(\epsilon^3)
\end{align}

When $x_i = 0$, false positives contribute $O(\epsilon^2)$ terms. Summing over all components:
\begin{align}
\E[\|\mathbf{x} - \hat{\mathbf{x}}\|^2] &= \sum_{i=1}^n p_i \E[(x_i - \hat{x}_i)^2 | x_i = 1] + O(\epsilon^2) \\
&= \sum_{i=1}^n p_i (1 - \|\mathbf{w}_i\|^2)^2 + \sum_{i=1}^n \sum_{j \neq i} p_i p_j \langle \mathbf{w}_i, \mathbf{w}_j \rangle^2 + O(\epsilon^3) \\
&= \mathcal{L}_{\text{cap}}(\mathbf{W}) + \mathcal{L}_{\text{int}}(\mathbf{W}, P) + O(\epsilon^3)
\end{align}
where $P_{ij} = p_i p_j$ for $i \neq j$ and $P_{ii} = p_i$.
\end{proof}

\subsection{Proof of Part (b): Reasoning Superposition}

\begin{lemma}
\label{lem:reason_expansion}
For continuous thought with trace embeddings $\mathbf{v}_k$ and attention weights $\alpha_k$, minimizing the next-token prediction loss (cross-entropy) implies minimizing the interference loss $\mathcal{L}_{\text{int}}$.
\end{lemma}

\begin{proof}
From Equation~\ref{eq:thought_decomp}, the thought vector is $\mathbf{h}_t = \sum_k \alpha_k^{(t)} \mathbf{v}_k$. The model predicts the next token by computing logits $\mathbf{z} = \mathbf{W}_U \mathbf{h}_T$ where $\mathbf{W}_U$ is the unembedding matrix.

\textbf{Step 1: Decomposing the logit.}
For the correct answer token $y^*$, let $\mathbf{u}_{y^*}$ be its unembedding vector. Suppose trace $k^*$ leads to the correct answer, i.e., $\mathbf{v}_{k^*}$ is aligned with $\mathbf{u}_{y^*}$. The logit for the correct answer is:
\begin{equation}
z_{y^*} = \langle \mathbf{u}_{y^*}, \mathbf{h}_T \rangle = \alpha_{k^*} \langle \mathbf{u}_{y^*}, \mathbf{v}_{k^*} \rangle + \underbrace{\sum_{j \neq k^*} \alpha_j \langle \mathbf{u}_{y^*}, \mathbf{v}_j \rangle}_{\text{interference noise}}
\end{equation}

\textbf{Step 2: Cross-entropy decomposition.}
The cross-entropy loss is $\mathcal{L}_{\text{CE}} = -\log \text{softmax}(z_{y^*})$. For small interference, Taylor expansion gives:
\begin{align}
\mathcal{L}_{\text{CE}} &\approx \mathcal{L}_{\text{ideal}} + \lambda \cdot \text{Var}[\text{interference}] + O(\epsilon^3) \\
&= \mathcal{L}_{\text{ideal}} + \lambda \sum_{j \neq k^*} \alpha_j^2 \langle \mathbf{u}_{y^*}, \mathbf{v}_j \rangle^2 + O(\epsilon^3)
\end{align}
where $\mathcal{L}_{\text{ideal}}$ is the loss with no interference and $\lambda > 0$ depends on the softmax temperature.

\textbf{Step 3: Connecting to $\mathcal{L}_{\text{int}}$.}
When traces $j$ and $k^*$ are both valid (both could lead to the answer), their embeddings must be distinguishable to avoid the model confusing them. Averaging over problems where different traces are correct:
\begin{equation}
\E[\mathcal{L}_{\text{CE}}] \approx \E[\mathcal{L}_{\text{ideal}}] + \lambda' \sum_{\substack{i, j \in V \\ i \neq j}} \langle \mathbf{v}_i, \mathbf{v}_j \rangle^2 + O(\epsilon^3)
\end{equation}

The second term is precisely $\mathcal{L}_{\text{int}}(\mathbf{W}, P)$ with $P_{ij} = \ind[\text{both valid}]$.

\textbf{Step 4: Capacity loss from normalization.}
For stable training, embeddings must have consistent norms. If $\|\mathbf{v}_k\| \ll 1$, the trace contributes negligibly to $\mathbf{h}_t$; if $\|\mathbf{v}_k\| \gg 1$, gradients become unstable. Layer normalization and weight decay implicitly enforce $\|\mathbf{v}_k\| \approx 1$, contributing the $\mathcal{L}_{\text{cap}}$ term.

Thus, minimizing $\mathcal{L}_{\text{CE}}$ implies minimizing $\mathcal{L}(\mathbf{W}, P) = \mathcal{L}_{\text{cap}} + \mathcal{L}_{\text{int}}$.
\end{proof}

\subsection{Completing the Proof of Theorem~\ref{thm:equivalence}}

\begin{proof}[Proof of Theorem~\ref{thm:equivalence}]
Part (a) follows from Lemma~\ref{lem:rep_expansion}: setting $P_{ij} = p_i p_j$ recovers the expected reconstruction error up to higher-order terms that vanish as sparsity increases.

Part (b) follows from Lemma~\ref{lem:reason_expansion}: setting $P_{ij} = \ind[\text{both valid}]$ recovers the implicit objective of the attention mechanism.

Both are instances of minimizing $\mathcal{L}(\mathbf{W}, P) = \mathcal{L}_{\text{cap}}(\mathbf{W}) + \mathcal{L}_{\text{int}}(\mathbf{W}, P)$, establishing the equivalence.
\end{proof}

\begin{remark}[Role of the Co-occurrence Matrix]
The co-occurrence matrix $P$ captures the \emph{structure} of the superposition problem:
\begin{itemize}
    \item \textbf{Representational:} $P_{ij} = p_i p_j$ reflects independent feature activations. High sparsity (small $p_i$) yields small $P_{ij}$, reducing interference cost.
    \item \textbf{Reasoning:} $P_{ij} = \ind[\text{both valid}]$ reflects the branching structure of valid traces. Problems with few valid paths have sparse $P$, enabling more superposition.
\end{itemize}
The unified framework reveals that both phenomena exploit the same mathematical structure: sparse co-occurrence enables packing more objects into limited dimensions.
\end{remark}

\begin{remark}[Interference as Noise vs.\ Exploration]
\label{rem:interference_exploration}
While the mathematics of superposition is equivalent across settings, the \emph{interpretation} of interference differs:

\paragraph{Representational setting: Interference is pure noise.} In feature encoding, interference $\langle \mathbf{w}_i, \mathbf{w}_j \rangle \neq 0$ is always undesirable---it causes feature $j$'s activation to corrupt the reconstruction of feature $i$. The goal is to minimize interference subject to capacity constraints.

\paragraph{Reasoning setting: Interference enables exploration.} In reasoning superposition, maintaining multiple traces simultaneously is a \emph{feature}, not a bug. The superposition $\mathbf{h}_t = \sum_k \alpha_k \mathbf{v}_k$ allows parallel exploration of the solution space, analogous to beam search or BFS. \citet{zhu2025reasoning} show this enables solving problems that sequential (DFS-like) reasoning cannot.

\paragraph{Resolution: Eventual collapse.} The key insight is that reasoning superposition must \emph{eventually} collapse to a single answer. The loss function $\mathcal{L}_{\text{int}}$ does not penalize superposition per se---it penalizes \emph{indistinguishability} between co-valid traces. When traces $i$ and $j$ are both valid, we need $\langle \mathbf{v}_i, \mathbf{v}_j \rangle \approx 0$ so that the model can later distinguish them and select the correct one. Thus:
\begin{itemize}
    \item \textbf{During exploration:} Multiple traces coexist in superposition, with attention weights $\alpha_k$ distributed across valid options.
    \item \textbf{At decision time:} The gating mechanism (Section~\ref{sec:timespace}) collapses the superposition, concentrating weight on the correct trace: $\alpha_{k^*} \to 1$.
\end{itemize}
Low interference $\langle \mathbf{v}_i, \mathbf{v}_j \rangle \approx 0$ ensures this collapse can be performed accurately---the traces remain distinguishable even while superposed. This parallels quantum computing, where orthogonal states can be maintained in superposition and later measured to yield a definite outcome.
\end{remark}

\section{Capacity Bounds}
\label{sec:capacity}

We now establish how many objects can be encoded in superposition as a function of dimension and sparsity.

\subsection{Sparsity Measure}

\begin{definition}[Co-occurrence Sparsity]
\label{def:sparsity}
For co-occurrence matrix $P \in [0,1]^{n \times n}$, define the \emph{sparsity} as:
\begin{equation}
S(P) = 1 - \frac{1}{n(n-1)} \sum_{i \neq j} P_{ij}
\end{equation}
This measures the fraction of pairs that rarely co-occur. $S = 1$ means no pairs co-occur; $S = 0$ means all pairs always co-occur.
\end{definition}

\begin{example}[Sparsity in Both Settings]
\label{ex:sparsity}
\begin{itemize}
    \item \textbf{Representational:} With $P_{ij} = p_i p_j$ and uniform $p_i = p$, we have $S = 1 - p^2$. For $p = 0.1$, sparsity is $S = 0.99$.
    \item \textbf{Reasoning:} For a binary tree with depth $D$, there are $2^D$ leaf traces, but only $O(D)$ pairs share any edge. Thus $S \approx 1 - O(D/2^D) \to 1$ as $D \to \infty$.
\end{itemize}
\end{example}

\begin{remark}[Global vs.\ Local Sparsity]
\label{rem:global_local}
A subtle distinction arises between \emph{global} and \emph{local} sparsity in reasoning:
\begin{itemize}
    \item \textbf{Global $P$:} The co-occurrence matrix over \emph{all possible} traces. This is sparse because most trace pairs never both lead to valid solutions.
    \item \textbf{Local $P^{(t)}$:} The co-occurrence matrix over traces \emph{active at timestep $t$}. During BFS-like exploration, the ``frontier'' of active traces can grow, making $P^{(t)}$ locally dense.
\end{itemize}

This apparent tension resolves as follows: even when the frontier is large, the \emph{relevant} traces (those contributing to the eventual answer) remain sparse. The model maintains superposition over the full frontier, but the interference that \emph{matters}---between traces that both reach the goal---remains controlled by global sparsity.

Formally, let $F_t$ denote the frontier at step $t$ with $|F_t| = O(b^t)$ for branching factor $b$. The local interference is:
\begin{equation}
\mathcal{L}_{\text{int}}^{(t)} = \sum_{i, j \in F_t} P^{(t)}_{ij} \langle \mathbf{v}_i, \mathbf{v}_j \rangle^2
\end{equation}
Even if $|F_t|$ is large, the capacity bound (Theorem~\ref{thm:capacity}) applies with dimension $d$: superposition succeeds when $|F_t| \leq O(d^2 \cdot S^{(t)} / \tau)$. The model handles frontier growth by either (a) having sufficient dimension $d$, or (b) using iterative refinement (Section~\ref{sec:timespace}) to prune the frontier over time.
\end{remark}

\subsection{Main Capacity Result}

\begin{theorem}[Capacity-Sparsity Bound]
\label{thm:capacity}
Let $\mathbf{W}^* = \arg\min_{\mathbf{W}} \mathcal{L}(\mathbf{W}, P)$ be an optimal encoding in $d$ dimensions. Define the \emph{effective number of encoded objects} as:
\begin{equation}
n_{\text{eff}} = |\{i : \|\mathbf{w}^*_i\| \geq 1/2\}|
\end{equation}
Then for any interference tolerance $\tau > 0$:\footnote{We use $\tau$ (tolerance) for the capacity constraint budget to distinguish from $\epsilon$ used elsewhere for approximation error bounds.}
\begin{equation}
n_{\text{eff}} \leq d + \frac{d(d-1)}{2} \cdot \frac{S(P)}{\tau}
\label{eq:capacity_bound}
\end{equation}
where $S(P)$ is the co-occurrence sparsity.
\end{theorem}

\begin{proof}
Let $\mathbf{G} = \mathbf{W}^\top \mathbf{W} \in \R^{n \times n}$ be the Gram matrix, with $G_{ij} = \langle \mathbf{w}_i, \mathbf{w}_j \rangle$.

\textbf{Step 1: Counting constraints.}
In $d$ dimensions, at most $d$ vectors can be mutually orthogonal. For $n_{\text{eff}} > d$ objects with $\|\mathbf{w}_i\| \geq 1/2$, there must exist pairs with $G_{ij} \neq 0$.

\textbf{Step 2: Interference budget.}
The interference loss is:
\begin{equation}
\mathcal{L}_{\text{int}} = \sum_{i \neq j} P_{ij} G_{ij}^2 \leq \tau
\end{equation}

By Cauchy-Schwarz, for unit-norm vectors $|G_{ij}| \leq 1$. For vectors with $\|\mathbf{w}_i\| \geq 1/2$, we have $|G_{ij}| \leq \|\mathbf{w}_i\| \|\mathbf{w}_j\| \leq 1$.

\textbf{Step 3: Packing argument.}
Consider the $\binom{n_{\text{eff}}}{2}$ pairs of effectively encoded objects. Each pair $(i,j)$ with $P_{ij} > 0$ contributes at least $P_{ij} G_{ij}^2$ to the interference loss.

The number of pairs with $P_{ij} > 0$ is at most $(1 - S) \cdot n_{\text{eff}}(n_{\text{eff}}-1)/2$.

For the interference constraint to be satisfied, we need:
\begin{equation}
\sum_{i \neq j} P_{ij} G_{ij}^2 \leq \tau
\end{equation}

\textbf{Step 4: Welch bound packing.}
Beyond the first $d$ orthogonal vectors, each additional vector must have non-zero inner product with at least one existing vector. The \emph{Welch bound} \citep{welch1974lower} provides a fundamental limit: for $n$ unit vectors in $\R^d$, the maximum coherence satisfies
\begin{equation}
\mu_{\max} := \max_{i \neq j} |G_{ij}| \geq \sqrt{\frac{n - d}{d(n-1)}}
\end{equation}
This implies that packing $n \gg d$ vectors requires $\mu_{\max} = \Omega(1/\sqrt{d})$, so the number of nearly-orthogonal directions (with $|G_{ij}| \leq \delta$) scales as $O(d^2/\delta^2)$.

Specifically, if we allow $|G_{ij}|^2 \leq \delta$ for pairs with $P_{ij} = 1-S$ (the ``dense'' pairs), the maximum packing is:
\begin{equation}
n_{\text{eff}} \leq d + \frac{d(d-1)}{2} \cdot \frac{1}{\delta}
\end{equation}

\textbf{Step 5: Combining with sparsity.}
The effective interference per dense pair is $P_{ij} G_{ij}^2 \leq (1-S) \cdot \delta$. To achieve total interference $\leq \tau$:
\begin{equation}
(1-S) \cdot \delta \cdot \frac{n_{\text{eff}}(n_{\text{eff}}-1)}{2} \leq \tau
\end{equation}

Solving for $n_{\text{eff}}$ with the packing constraint yields:
\begin{equation}
n_{\text{eff}} \leq d + \frac{d(d-1)}{2} \cdot \frac{S}{\tau}
\end{equation}
as claimed.
\end{proof}

\begin{corollary}[Linear Scaling with Sparsity]
\label{cor:linear}
For fixed dimension $d$ and interference tolerance $\tau$, the capacity scales linearly with sparsity:
\begin{equation}
n_{\text{eff}} = O(d^2 \cdot S / \tau)
\end{equation}
\end{corollary}

\begin{corollary}[Equivalence of Scaling]
\label{cor:equiv_scaling}
Both representational and reasoning superposition exhibit the same capacity scaling:
\begin{itemize}
    \item \textbf{Representational:} With activation probability $p$, sparsity is $S = 1 - p^2$, giving $n_{\text{eff}} = O(d^2 / p^2)$ for small $p$.
    \item \textbf{Reasoning:} With $K$ total traces and $k$ valid traces, if validity is sparse ($k \ll K$), then $S \approx 1$ and $n_{\text{eff}} = O(d^2)$.
\end{itemize}
The identical scaling confirms the phenomena are mathematically equivalent.
\end{corollary}

\subsection{Tightness of the Bound}

\begin{proposition}[Lower Bound]
\label{prop:lower}
There exist co-occurrence matrices $P$ and encodings $\mathbf{W}$ achieving:
\begin{equation}
n_{\text{eff}} \geq d + c \cdot d^2 \cdot S / \tau
\end{equation}
for a universal constant $c > 0$.
\end{proposition}

\begin{proof}[Proof Sketch]
Consider the construction of \citet{elhage2022toy} using regular polytopes. For $d = 2$ dimensions, $n = 2k$ features can be arranged as vertices of a regular $k$-gon and its reflection, achieving $\langle \mathbf{w}_i, \mathbf{w}_j \rangle = \cos(2\pi |i-j|/k)$.

For sparse features with $P_{ij} = p^2$, the interference is:
\begin{equation}
\mathcal{L}_{\text{int}} = p^2 \sum_{i \neq j} \cos^2(2\pi|i-j|/k) = O(p^2 k)
\end{equation}

Setting this equal to $\tau$ and solving: $k = O(\tau / p^2) = O(\tau \cdot d^2 \cdot S)$ for $d = 2$.

The construction generalizes to higher dimensions using cross-polytopes and their duals.
\end{proof}

\section{Time-Space Equivalence}
\label{sec:timespace}

A key architectural question is how iterative refinement (temporal depth) relates to representational capacity (spatial width). We prove that time can compensate for space in superposition.

\subsection{Setup}

Consider a model processing information through a bottleneck of dimension $d$, with $T$ iterative refinement steps. Let $f_\theta: \R^d \to \R^d$ denote the transformation at each step, and let $\mathbf{h}_0 \in \R^d$ be the initial (superposed) representation.

\begin{definition}[Iterative Refinement Model]
\label{def:iterative}
An \emph{iterative refinement model} $\mathcal{M}_T(d)$ computes:
\begin{equation}
\mathbf{h}_t = f_\theta(\mathbf{h}_{t-1}), \quad t = 1, \ldots, T
\end{equation}
The final output is decoded from $\mathbf{h}_T$.
\end{definition}

\begin{definition}[Effective Capacity]
\label{def:eff_cap}
The \emph{effective capacity} $\mathcal{I}_T(d)$ of model $\mathcal{M}_T(d)$ is the maximum number of objects that can be encoded and recovered with error at most $\epsilon$.
\end{definition}

\subsection{Main Time-Space Result}

\begin{theorem}[Time-Space Tradeoff]
\label{thm:timespace}
Suppose $f_\theta$ includes a gating mechanism capable of computing:
\begin{equation}
f_\theta\left(\sum_k \alpha_k \mathbf{v}_k\right) = \sum_k \alpha_k \cdot g_k(\boldsymbol{\alpha}, \mathbf{V}) \cdot \mathbf{v}_k
\label{eq:gating}
\end{equation}
where $g_k: \R^K \times \R^{d \times K} \to \R$ are learnable gating functions. Then:
\begin{equation}
\mathcal{I}_T(d) \geq \mathcal{I}_1(d \cdot 2^{T-1})
\end{equation}
That is, $T$ iterative steps provide capacity equivalent to exponentially more dimensions.
\end{theorem}

\begin{proof}
We prove by induction that $\mathcal{M}_T(d)$ can simulate $\mathcal{M}_1(d \cdot 2^{T-1})$.

\textbf{Base case ($T = 1$):} Trivially, $\mathcal{I}_1(d) = \mathcal{I}_1(d \cdot 2^0)$.

\textbf{Inductive step:} Assume $\mathcal{M}_T(d)$ can simulate $\mathcal{M}_1(d \cdot 2^{T-1})$. We show $\mathcal{M}_{T+1}(d)$ can simulate $\mathcal{M}_1(d \cdot 2^T)$.

Consider $K = 2^T$ objects encoded in $d \cdot 2^T$ dimensions. Partition them into two groups $G_0, G_1$ of size $2^{T-1}$ each, where group $G_b$ uses dimensions in the range $[b \cdot d \cdot 2^{T-1}, (b+1) \cdot d \cdot 2^{T-1})$.

The key observation is that the gating mechanism (Equation~\ref{eq:gating}) can implement \emph{conditional computation}:
\begin{equation}
g_k(\boldsymbol{\alpha}, \mathbf{V}) = \begin{cases}
1 & \text{if } k \in G_{b^*} \\
0 & \text{otherwise}
\end{cases}
\end{equation}
where $b^* \in \{0, 1\}$ is determined by the context (e.g., which group is more relevant to the query).

After the first step of $\mathcal{M}_{T+1}(d)$:
\begin{equation}
\mathbf{h}_1 = \sum_{k \in G_{b^*}} \alpha_k \mathbf{v}_k
\end{equation}

This effectively ``selects'' one group, reducing the problem to encoding $2^{T-1}$ objects. By the inductive hypothesis, the remaining $T$ steps can handle this with capacity equivalent to $d \cdot 2^{T-1}$ dimensions.

Since the first step doubled the effective dimension (by selecting one of two groups), and the remaining steps provide $d \cdot 2^{T-1}$ capacity:
\begin{equation}
\mathcal{I}_{T+1}(d) \geq 2 \cdot d \cdot 2^{T-1} = d \cdot 2^T = \mathcal{I}_1(d \cdot 2^T)
\end{equation}
completing the induction.
\end{proof}

\subsection{Interpretation}

\begin{remark}[Unpacking Superposition]
The proof reveals the mechanism by which iterative refinement ``unpacks'' superposition:
\begin{enumerate}
    \item Initially, $\mathbf{h}_0$ contains all $K$ objects in superposition.
    \item Each step $t$ uses gating to selectively amplify relevant objects and suppress irrelevant ones.
    \item After $T$ steps, only the target object remains, enabling accurate decoding.
\end{enumerate}
This is analogous to binary search: each step halves the candidate set, achieving exponential speedup.
\end{remark}

\begin{remark}[Connection to Reasoning]
In reasoning superposition, this corresponds to:
\begin{enumerate}
    \item $\mathbf{h}_0$ contains all valid reasoning traces in superposition.
    \item Each thought step refines the weights $\alpha_k$, increasing weight on promising traces.
    \item The final step ``collapses'' to the correct answer, analogous to measurement in quantum computing.
\end{enumerate}
\end{remark}

\subsection{Lower Bound}

\begin{theorem}[Time-Space Lower Bound]
\label{thm:timespace_lower}
For any model $\mathcal{M}_T(d)$ without the gating mechanism (i.e., $f_\theta$ is a fixed linear transformation plus nonlinearity):
\begin{equation}
\mathcal{I}_T(d) \leq \mathcal{I}_1(d \cdot T)
\end{equation}
That is, without gating, time provides at most linear (not exponential) capacity increase.
\end{theorem}

\begin{proof}[Proof Sketch]
Without gating, each step applies a fixed transformation. The composition of $T$ linear transformations is still linear (rank at most $d$). Nonlinearities can increase expressiveness, but for reconstruction tasks, the information bottleneck limits capacity to $O(d \cdot T)$ bits.

Formally, consider the information flow: $\mathbf{h}_0 \to \mathbf{h}_1 \to \cdots \to \mathbf{h}_T$. Each $\mathbf{h}_t \in \R^d$ can carry at most $O(d)$ bits of information. Even with $T$ steps, the total information is at most $O(d \cdot T)$, giving capacity $\mathcal{I}_T(d) \leq O(d \cdot T)$.
\end{proof}

\begin{corollary}[Necessity of Gating]
\label{cor:gating}
The exponential time-space tradeoff (Theorem~\ref{thm:timespace}) requires the gating mechanism. This explains why architectures like Transformers (with attention-based gating) and LSTMs (with input/forget gates) can achieve greater effective capacity than feedforward networks of the same width.
\end{corollary}

\subsection{Architectural Implications}

Our results suggest a design principle for models combining representational and reasoning superposition:

\begin{enumerate}
    \item \textbf{Bottleneck layer:} Compress representations to dimension $d \ll D$ (where $D$ is the original dimension), forcing representational superposition.

    \item \textbf{Iterative refinement:} Add $T$ steps of gated transformation in the bottleneck space, enabling reasoning superposition to unpack the compressed representations.

    \item \textbf{Capacity planning:} By Theorem~\ref{thm:timespace}, the effective capacity is $\mathcal{I}_T(d) \geq \mathcal{I}_1(d \cdot 2^{T-1})$. Choose $d$ and $T$ such that $d \cdot 2^{T-1} \geq$ (number of features/traces to encode).
\end{enumerate}

\begin{example}[Concrete Design]
To encode $n = 1000$ features with a bottleneck of $d = 64$ dimensions:
\begin{itemize}
    \item Single-step model: Can encode $\mathcal{I}_1(64) \approx 64$ features orthogonally, or $O(64^2) \approx 4000$ with superposition (if sparse).
    \item With $T = 4$ steps: Effective capacity $\mathcal{I}_4(64) \geq \mathcal{I}_1(64 \cdot 8) = \mathcal{I}_1(512)$, sufficient for 1000 features even without superposition.
\end{itemize}
\end{example}

\section{Discussion}
\label{sec:discussion}

\subsection{Summary of Results}

We have established three main results:

\begin{enumerate}
    \item \textbf{Equivalence (Theorem~\ref{thm:equivalence}):} Representational and reasoning superposition minimize the same loss function $\mathcal{L}(\mathbf{W}, P) = \mathcal{L}_{\text{cap}} + \mathcal{L}_{\text{int}}$, differing only in the co-occurrence matrix $P$.

    \item \textbf{Capacity (Theorem~\ref{thm:capacity}):} Both phenomena exhibit identical scaling: $n_{\text{eff}} = O(m^2 \cdot S / \epsilon)$, where sparsity $S$ enables encoding more objects.

    \item \textbf{Time-Space (Theorem~\ref{thm:timespace}):} Iterative refinement with gating provides exponential capacity increase: $\mathcal{I}_T(d) \geq \mathcal{I}_1(d \cdot 2^{T-1})$.
\end{enumerate}

\subsection{Implications for Mechanistic Interpretability}

Our results suggest several directions for interpretability research:

\paragraph{Transfer of techniques.} Methods developed for representational superposition may apply to reasoning superposition:
\begin{itemize}
    \item \textbf{Sparse autoencoders} \citep{bricken2023monosemanticity} disentangle superposed features by learning an overcomplete basis. Analogous methods could disentangle superposed reasoning traces.
    \item \textbf{Feature steering} \citep{templeton2024scaling} modifies activations to change model behavior. Similar interventions on thought vectors could steer reasoning.
\end{itemize}

\paragraph{Unified analysis framework.} Rather than analyzing features and reasoning traces separately, our framework suggests analyzing the joint space of ``objects in superposition'' with their co-occurrence structure $P$.

\paragraph{Predicting polysemanticity.} The capacity bounds (Theorem~\ref{thm:capacity}) predict when neurons will be polysemantic: when $n > m + m^2 S / \epsilon$, some neurons must encode multiple features. This could guide architecture design for interpretability.

\subsection{Implications for Architecture Design}

\paragraph{Width vs. depth tradeoff.} Theorem~\ref{thm:timespace} quantifies a fundamental tradeoff:
\begin{itemize}
    \item \textbf{Wider models} ($d$ large): More capacity per step, less superposition needed.
    \item \textbf{Deeper/iterative models} ($T$ large): Exponential capacity from iteration, more superposition tolerated.
\end{itemize}
For interpretability, wider models may be preferable (less superposition to disentangle), but iterative models may be more parameter-efficient.

\paragraph{Bottleneck design.} Our results suggest deliberate bottlenecks can be useful:
\begin{enumerate}
    \item Force superposition by constraining $d < n$.
    \item Add iterative steps to unpack superposition during inference.
    \item The bottleneck acts as an ``interpretability checkpoint'' where all information must pass through a low-dimensional space.
\end{enumerate}

\subsection{Implications for AI Safety}

\paragraph{Understanding reasoning.} If language models reason via superposition of traces (as suggested by \citet{zhu2025reasoning}), then:
\begin{itemize}
    \item Harmful reasoning may be one of multiple superposed traces.
    \item Steering away from harm requires identifying and suppressing specific traces without disrupting beneficial ones.
    \item The co-occurrence structure $P$ determines which traces can be independently manipulated.
\end{itemize}

\paragraph{Predictable failures.} The capacity bounds predict when superposition-induced errors occur: when the number of active objects exceeds $m + m^2 S / \epsilon$, interference causes errors. This could enable proactive safety measures.

\subsection{Limitations and Future Work}

\paragraph{Approximations in proofs.} Our equivalence theorem (Theorem~\ref{thm:equivalence}) relies on a small-interference approximation. For highly non-sparse settings, higher-order terms may matter.

\paragraph{Empirical validation.} While our results are theoretical, empirical validation would strengthen the claims:
\begin{itemize}
    \item Measure whether the capacity scaling $O(m^2 S / \epsilon)$ holds in practice.
    \item Test whether iterative models achieve the predicted exponential capacity increase.
    \item Verify that techniques transfer between representational and reasoning settings.
\end{itemize}

\paragraph{Beyond binary co-occurrence.} Our framework uses binary validity for reasoning traces. Extending to graded validity (partial correctness) could better capture real reasoning.

\paragraph{Learning dynamics.} We characterize optimal solutions but not learning dynamics. Understanding \emph{how} superposition emerges during training \citep{zhu2025emergence} remains important.

\subsection{Conclusion}

We have proven that representational superposition and reasoning superposition are mathematically equivalent, both arising from optimizing capacity subject to interference constraints weighted by co-occurrence. This unification suggests that the extensive literature on feature superposition directly informs our understanding of reasoning in neural networks, and vice versa. The time-space equivalence theorem further reveals that iterative computation can exponentially amplify representational capacity, providing theoretical foundations for architectures that combine both mechanisms.


\bibliography{references}
\bibliographystyle{plain}

\appendix
\section{Extended Proofs}
\label{app:proofs}

\subsection{Proof of Lemma~\ref{lem:rep_expansion} (Full Details)}

We provide a complete derivation of the expected reconstruction error.

\begin{proof}
Let $\mathbf{x} \in \{0,1\}^n$ with independent components $x_i \sim \text{Bernoulli}(p_i)$. The pre-activation output is:
\begin{equation}
\mathbf{z} = \mathbf{W}^\top \mathbf{W} \mathbf{x}
\end{equation}
with components:
\begin{equation}
z_i = \sum_{j=1}^n x_j \langle \mathbf{w}_i, \mathbf{w}_j \rangle = x_i \|\mathbf{w}_i\|^2 + \sum_{j \neq i} x_j \langle \mathbf{w}_i, \mathbf{w}_j \rangle
\end{equation}

After ReLU: $\hat{x}_i = \max(0, z_i)$.

\textbf{Case 1: $x_i = 1$.}
\begin{equation}
z_i = \|\mathbf{w}_i\|^2 + \sum_{j \neq i} x_j \langle \mathbf{w}_i, \mathbf{w}_j \rangle
\end{equation}

Define $\delta_i = \sum_{j \neq i} x_j \langle \mathbf{w}_i, \mathbf{w}_j \rangle$ (the interference term).

For $\|\mathbf{w}_i\|^2 + \delta_i > 0$ (which holds when $\delta_i > -\|\mathbf{w}_i\|^2$):
\begin{equation}
\hat{x}_i = \|\mathbf{w}_i\|^2 + \delta_i
\end{equation}

The squared error is:
\begin{align}
(x_i - \hat{x}_i)^2 &= (1 - \|\mathbf{w}_i\|^2 - \delta_i)^2 \\
&= (1 - \|\mathbf{w}_i\|^2)^2 - 2(1 - \|\mathbf{w}_i\|^2)\delta_i + \delta_i^2
\end{align}

Taking expectation over $\{x_j\}_{j \neq i}$:
\begin{align}
\E[\delta_i] &= \sum_{j \neq i} p_j \langle \mathbf{w}_i, \mathbf{w}_j \rangle \\
\E[\delta_i^2] &= \sum_{j \neq i} p_j \langle \mathbf{w}_i, \mathbf{w}_j \rangle^2 + \sum_{j \neq k, j,k \neq i} p_j p_k \langle \mathbf{w}_i, \mathbf{w}_j \rangle \langle \mathbf{w}_i, \mathbf{w}_k \rangle
\end{align}

For sparse features, the cross-terms are $O(p^2)$ and can be absorbed into higher-order corrections:
\begin{equation}
\E[(x_i - \hat{x}_i)^2 | x_i = 1] = (1 - \|\mathbf{w}_i\|^2)^2 + \sum_{j \neq i} p_j \langle \mathbf{w}_i, \mathbf{w}_j \rangle^2 + O(\epsilon^3)
\end{equation}

\textbf{Case 2: $x_i = 0$.}

False positive: $\hat{x}_i > 0$ when $\sum_{j \neq i} x_j \langle \mathbf{w}_i, \mathbf{w}_j \rangle > 0$.

The expected squared error from false positives is:
\begin{equation}
\E[(0 - \hat{x}_i)^2 | x_i = 0] = \E[\delta_i^2 \cdot \ind[\delta_i > 0]] = O(\epsilon^2)
\end{equation}
for sparse features.

\textbf{Combining cases:}
\begin{align}
\E[(x_i - \hat{x}_i)^2] &= p_i \E[(x_i - \hat{x}_i)^2 | x_i = 1] + (1-p_i) \E[(x_i - \hat{x}_i)^2 | x_i = 0] \\
&= p_i \left[(1 - \|\mathbf{w}_i\|^2)^2 + \sum_{j \neq i} p_j \langle \mathbf{w}_i, \mathbf{w}_j \rangle^2\right] + O(\epsilon^2)
\end{align}

Summing over all $i$:
\begin{equation}
\E[\|\mathbf{x} - \hat{\mathbf{x}}\|^2] = \sum_i p_i (1 - \|\mathbf{w}_i\|^2)^2 + \sum_{i \neq j} p_i p_j \langle \mathbf{w}_i, \mathbf{w}_j \rangle^2 + O(\epsilon^3)
\end{equation}

This equals $\mathcal{L}(\mathbf{W}, P)$ with $P_{ij} = p_i p_j$.
\end{proof}

\subsection{Tighter Capacity Bounds}

\begin{proposition}[Refined Capacity Bound]
\label{prop:tight_capacity}
Under additional assumptions on the co-occurrence structure, the capacity bound can be tightened to:
\begin{equation}
n_{\text{eff}} \leq m + \frac{m^2}{2} \cdot \min\left(\frac{S}{\epsilon}, \frac{1}{\sqrt{\epsilon}}\right)
\end{equation}
\end{proposition}

\begin{proof}[Proof Sketch]
The refinement comes from noting that for very small $\epsilon$, the constraint $P_{ij} G_{ij}^2 \leq \epsilon$ for all pairs forces $|G_{ij}| \leq \sqrt{\epsilon / P_{ij}}$. When $P_{ij} = \Theta(1)$ for some pairs, this gives $|G_{ij}| \leq \sqrt{\epsilon}$, limiting how many vectors can pack into the space regardless of sparsity.
\end{proof}

\section{Additional Examples}
\label{app:examples}

\subsection{Polytope Constructions}

\begin{example}[Pentagon in 2D]
For $m = 2$ and $n = 5$, the optimal construction places features at vertices of a regular pentagon:
\begin{equation}
\mathbf{w}_k = \begin{pmatrix} \cos(2\pi k / 5) \\ \sin(2\pi k / 5) \end{pmatrix}, \quad k = 0, \ldots, 4
\end{equation}
The interference between adjacent features is $\langle \mathbf{w}_k, \mathbf{w}_{k+1} \rangle = \cos(2\pi/5) \approx 0.309$.

For uniform sparsity $p$, the total interference is:
\begin{equation}
\mathcal{L}_{\text{int}} = p^2 \cdot 5 \cdot 4 \cdot \cos^2(2\pi/5) / 2 \approx 1.9 p^2
\end{equation}
\end{example}

\subsection{Reasoning Trace Example}

\begin{example}[Graph Reachability]
Consider a directed graph with source $s$ and target $t$, with three valid paths:
\begin{itemize}
    \item Path 1: $s \to a \to t$
    \item Path 2: $s \to b \to t$
    \item Path 3: $s \to a \to b \to t$
\end{itemize}

The co-validity matrix is:
\begin{equation}
P = \begin{pmatrix} 1 & 1 & 1 \\ 1 & 1 & 1 \\ 1 & 1 & 1 \end{pmatrix}
\end{equation}
(all paths are valid, so all pairs co-occur).

Sparsity: $S = 1 - 6/6 = 0$ (no sparsity, all pairs co-valid).

By Theorem~\ref{thm:capacity}, with $S = 0$, we need $n_{\text{eff}} \leq m$ --- no superposition benefit. The three traces must be nearly orthogonal, requiring $m \geq 3$ dimensions.

In contrast, if only paths 1 and 2 were valid (path 3 invalid), then $S = 1 - 2/6 = 2/3$, allowing more superposition.
\end{example}


\end{document}
